% Arquivo de configuracao do documento (template de artigo do IMC/UNIFEI)
% Define os dados dos autores e do trabalho. Itens marcados com \pendente{} ficam para a equipe.

\autorNomeA{João Leão}
\autorSobrenomeA{Guedes Filho}
\autorNomeB{Rafael Curiel}
\autorSobrenomeB{Monteiro}

% E-mails institucionais: a preencher pela equipe
\emailAutorA{[e-mail a preencher]}
\emailAutorB{[e-mail a preencher]}

\titulo{Detecção de Pessoas na Borda: Viabilidade da Inferência Local em Câmeras de Smart Home}

% Trabalho da disciplina XRSC10 (Tópicos em Redes e Sistemas Computacionais), sem orientador de TCC:
% o campo de orientação traz o professor da disciplina.
\orientadorNome{Prof. Dr. Laércio}
\orientadorSobrenome{Pioli}
\coorientadorNome{}
\coorientadorSobrenome{}
\txtOrientador{Professor}
\txtCoorientador{Coorientador}

\ano{2026}
\mes{Dezembro}
\local{Itajubá}

\palavraChaveA{Computação de borda}
\palavraChaveB{Detecção de pessoas}
\palavraChaveC{Quantização INT8}
\palavraChaveD{Casa inteligente}
\palavraChaveE{}

% A classe gera folha de aprovação (formato de TCC). Para a disciplina, só o professor.
\membroBancaA{Prof. Dr. Laércio Pioli}
\membroBancaB{}
\membroBancaC{}
\membroBancaD{}
\membroBancaE{}
\dataDefesa{[data da apresentação a preencher]}
\pdfAprovacao{}

% Textos da capa: a classe usa "Trabalho de Conclusão de Curso"; aqui é um artigo de disciplina.
\renewcommand{\tipo}{Artigo da disciplina XRSC10 -- Tópicos em Redes e Sistemas Computacionais}

% Pacotes e comandos adicionais deste trabalho
\usepackage{booktabs}
% Marcador visível para o que a equipe ainda precisa preencher
\newcommand{\pendente}[1]{\textcolor{red}{\textbf{[#1]}}}
% Números citados no texto: gerados por src/analysis/artigo_tex.py a partir dos dados (não edite)
% Gerado por src/analysis/artigo_tex.py a partir de results/. Não edite à mão.
\newcommand{\bandaQuatroCinco}{1,7}
\newcommand{\bandaQuatroQuinze}{5,2}
\newcommand{\bandaQuatroTrinta}{10,5}
\newcommand{\bandaSeisCinco}{3,0}
\newcommand{\bandaSeisQuinze}{8,9}
\newcommand{\bandaSeisTrinta}{17,7}
\newcommand{\bandaTresCinco}{0,8}
\newcommand{\bandaTresQuinze}{2,5}
\newcommand{\bandaTresTrinta}{5,0}
\newcommand{\etoeSsdFpQuatroNuvem}{102}
\newcommand{\etoeSsdFpTresNuvem}{87}
\newcommand{\etoeYoloFpQuatroNuvem}{111}
\newcommand{\etoeYoloFpSeisNuvem}{165}
\newcommand{\etoeYoloFpTresNuvem}{77}
\newcommand{\mapSsdFpQuatroNuvem}{0,595}
\newcommand{\mapSsdFpTresNuvem}{0,592}
\newcommand{\mapYoloFpQuatroNuvem}{0,745}
\newcommand{\mapYoloFpSeisNuvem}{0,761}
\newcommand{\mapYoloFpTresNuvem}{0,668}
\newcommand{\nCenarios}{15}
\newcommand{\nCombBorda}{0}
\newcommand{\nCombNuvem}{5}
\newcommand{\nQuadrosValidos}{300}

% Figuras geradas pela análise (e6)
\graphicspath{{./figs/}}
